\section{错版 Fixed vs Corrected Fixed}

\subsection{Bug 图}

\begin{figure}[H]
\centering
\includegraphics[width=0.92\textwidth]{fixed-hint-bug-task-depth-distribution.png}
\caption{old fixed 的 bug 会把不同任务的 hint depth 分布系统性压歪。}
\end{figure}

\subsection{任务级摘要}

\begin{table}[H]
\centering
\small
\setlength{\tabcolsep}{4pt}
\begin{tabular}{p{3.2cm} c c c p{4.1cm}}
\toprule
Task & Changed rate & Avg depth correct & Avg depth buggy & 主效应 \\
\midrule
\texttt{task1\_sid\_sft} & 7.64\% & 1.048 & 0.938 & 只被轻微压浅 \\
\texttt{task4\_hisTitle2sid} & 82.35\% & 1.533 & 0.188 & hardest task 几乎被压成 depth-0 \\
\texttt{task5\_title\_desc2sid} & 19.41\% & 0.519 & 0.245 & 一部分 depth-2 被压回 0/1 \\
\bottomrule
\end{tabular}
\caption{old fixed 的 bug 不是“hint 更准”，而是把任务分布系统性压浅。}
\end{table}

\subsection{结论}

\begin{itemize}
\item old fixed 的表面优势，不是来自更正确的 hint map，而是因为 hardest task 被系统性少给 hint。
\item 因此它更像一个带偏置的历史现象学上界，而不是应该继续复用的方法定义。
\item 这也是为什么本文后文在需要 fixed reference 时，默认优先引用 corrected \texttt{fixed taskfix sid-only}。
\end{itemize}
